

\section{Wrapping Convolution}

The nonwrapping convolution can enter the all-zero state after a run of
$\sigma$ zero outputs.  A sparse input can then leave the convolution inactive
for a long interval.  Expand--Convolute addresses this event by fixing the tap
on the oldest state bit to one \cite{EC}.  We call the resulting construction
the \emph{wrapping convolution}.  With zero initial state, its recurrence is
\begin{equation}
 y_i=x_i+y_{i-\sigma}
     +\sum_{j=1}^{\sigma-1}b_{i,j}y_{i-j},
 \qquad b_{i,j}\gets\{0,1\}\text{ uniformly},
 \label{eq:wrapping-recurrence}
\end{equation}
where indices at most zero contain zero.  All arithmetic is binary, and the
random taps are independent across positions.  This definition identifies the
fixed tap by the state bit it multiplies and therefore does not depend on a tap
ordering convention.

\subsection{Exact transfer enumerator}
\label{sec:wrapping-transfer-enumerator}

The same reduced state used for the nonwrapping convolution remains exact.
State $\ell<\sigma$ means that the output prefix ends in exactly $\ell$ zeros,
and state $\sigma$ means that its last $\sigma$ outputs are zero.  The only new
case is $\ell=\sigma-1$.  The feedback state is then
$10^{\sigma-1}$ in chronological order.  The fixed oldest tap makes the next
output equal to $1+x_i$.

Let $\mathbf W_\sigma(X,Y)$ be the $(\sigma+1)\times(\sigma+1)$ matrix whose
rows and columns use the state order $0,\ldots,\sigma$.  Its nonzero entries are
\begin{align}
 (\mathbf W_\sigma)_{\ell,0}
   &=\frac{1+X}{2}Y,
 &
 (\mathbf W_\sigma)_{\ell,\ell+1}
   &=\frac{1+X}{2}
 &&(0\leq\ell<\sigma-1),
 \label{eq:wrapping-bivariate-active}\\
 (\mathbf W_\sigma)_{\sigma-1,0}
   &=Y,
 &
 (\mathbf W_\sigma)_{\sigma-1,\sigma}
   &=X,
 \label{eq:wrapping-bivariate-critical}\\
 (\mathbf W_\sigma)_{\sigma,0}
   &=XY,
 &
 (\mathbf W_\sigma)_{\sigma,\sigma}
   &=1.
 \label{eq:wrapping-bivariate-zero}
\end{align}
The exponent of $X$ records input weight, and the exponent of $Y$ records
output weight.

\begin{lemma}[Exact wrapping convolution enumerator]
\label{lem:wrapping-transfer-enumerator}
Let $\mathbf e_\sigma$ be the unit column vector for state $\sigma$, and let
$\mathbf 1$ be the all-one column vector.  For every $u,h\in[0,n]$,
\begin{equation}
 \boxed{
 \enum^{\textsf{wConv}}_{u,h}
   =[X^uY^h],\mathbf e_\sigma^\mathsf{T}
     \mathbf W_\sigma(X,Y)^n\mathbf 1.
 }
 \label{eq:wrapping-transfer-enumerator}
\end{equation}
\end{lemma}

\begin{proof}
Suppose first that $\ell<\sigma-1$.  At least one of the $\sigma-1$ state
bits with a random tap is nonzero.  Their inner product is therefore uniform,
so the next output is uniform for either input bit.  This gives
\eqref{eq:wrapping-bivariate-active}.  In state $\sigma-1$, the fixed tap is
the only active tap.  Input zero produces output one and returns to state zero;
input one produces output zero and enters state $\sigma$.  This gives
\eqref{eq:wrapping-bivariate-critical}.  In state $\sigma$, the feedback is
zero and the output equals the input, which gives
\eqref{eq:wrapping-bivariate-zero}.  Matrix multiplication concatenates these
transitions.  The boundary vectors select the zero initial state and sum over
the final state.  Coefficient extraction proves the claim.
\end{proof}

For a Bernoulli input with nonzero probability $q$, define
$\mathbf N_{\sigma,q}(Y)$ by
\begin{align}
 (\mathbf N_{\sigma,q})_{\ell,0}
   &=\frac{Y}{2},
 &
 (\mathbf N_{\sigma,q})_{\ell,\ell+1}
   &=\frac{1}{2}
 &&(0\leq\ell<\sigma-1),
 \label{eq:wrapping-message-active}\\
 (\mathbf N_{\sigma,q})_{\sigma-1,0}
   &=(1-q)Y,
 &
 (\mathbf N_{\sigma,q})_{\sigma-1,\sigma}
   &=q,
 \label{eq:wrapping-message-critical}\\
 (\mathbf N_{\sigma,q})_{\sigma,0}
   &=qY,
 &
 (\mathbf N_{\sigma,q})_{\sigma,\sigma}
   &=1-q.
 \label{eq:wrapping-message-zero}
\end{align}

\begin{theorem}[Exact wrapped Expand--Convolute first moment]
\label{thm:wrapped-expand-convolute-first-moment}
Sample $\B\gets\dist^{\mathrm{Ber}}_{k,n,p}$ and an independent wrapping
convolution $\C$ of memory $\sigma$.  Set $\G:=\B\C$.  For $L\in[0,n]$, let
\begin{equation}
 \mu_{\mathrm{wEC}}(L)
   :=\sum_{r=1}^{k}\binom{k}{r}
     \sum_{h=0}^{L}[Y^h]\,
     \mathbf e_\sigma^\mathsf{T}
     \mathbf N_{\sigma,q_r}(Y)^n\mathbf 1.
 \label{eq:wrapped-expand-convolute-first-moment}
\end{equation}
Then
\begin{equation}
 \Pr\left[\exists\xv\neq0:\ \hw{\xv\G}\leq L\right]
 \leq \mu_{\mathrm{wEC}}(L).
 \label{eq:wrapped-expand-convolute-failure}
\end{equation}
\end{theorem}

\begin{proof}
Fix a message of weight $r$.  Its Bernoulli-expander image has independent
coordinates with nonzero probability $q_r$ by
Lemma~\ref{lem:bernoulli-expander-enumerator}.  Averaging the three transition
cases in Lemma~\ref{lem:wrapping-transfer-enumerator} over the input bit gives
$\mathbf N_{\sigma,q_r}$.  Its low-degree coefficients give the probability
that this message has output weight at most $L$.  Sum over nonzero messages
and apply Markov's inequality.
\end{proof}

For every $z\in(0,1]$, nonnegativity of the coefficients gives the computable
bound
\begin{equation}
 \sum_{h=0}^{L}[Y^h]\,
 \mathbf e_\sigma^\mathsf{T}\mathbf N_{\sigma,q}(Y)^n\mathbf 1
 \leq z^{-L}\mathbf e_\sigma^\mathsf{T}
 \mathbf N_{\sigma,q}(z)^n\mathbf 1.
 \label{eq:wrapped-expand-convolute-chernoff}
\end{equation}
The matrix has only $\sigma+1$ states, including for the wrapping construction.
Thus the exact enumerator avoids both a $2^\sigma$ full-state calculation and
the comparison chain used in the earlier analysis.

\subsection{Asymptotic distance}
\label{sec:wrapped-ec-asymptotic}

The transfer matrix also gives a closed distance theorem.  Let
$\rho_\sigma(q,z)$ denote the spectral radius of
$\mathbf N_{\sigma,q}(z)$.  For $q\in(0,1/2]$, define
\begin{equation}
 J^{\mathrm{w}}_{\sigma,\delta}(q)
 :=\sup_{0<z\leq1}
 \left\{\delta\ln z-\ln\rho_\sigma(q,z)\right\}.
 \label{eq:wrapped-ec-tail-rate}
\end{equation}
The sparse limit depends on
\begin{equation}
 I^{\mathrm{w}}_{\sigma,\delta}
 :=
 \left(
  \sqrt{1-2\delta+2^{1-\sigma}\delta}
  -\sqrt{2^{1-\sigma}\delta}
 \right)^2.
 \label{eq:wrapped-ec-sparse-rate}
\end{equation}

\begin{lemma}[Wrapped-convolution endpoint exponents]
\label{lem:wrapped-ec-endpoint-exponents}
Fix an integer $\sigma\geq1$ and $\delta\in(0,1/2)$.  The function
$J^{\mathrm{w}}_{\sigma,\delta}$ is continuous and strictly positive on
$(0,1/2]$.  Its endpoint exponents satisfy
\begin{align}
 \liminf_{q\downarrow0}
 \frac{J^{\mathrm{w}}_{\sigma,\delta}(q)}{q}
 &\geq I^{\mathrm{w}}_{\sigma,\delta},
 \label{eq:wrapped-ec-sparse-exponent-limit}\\
 J^{\mathrm{w}}_{\sigma,\delta}(1/2)
 &=\ln2-\mathsf{h}(\delta).
 \label{eq:wrapped-ec-dense-exponent}
\end{align}

Let $H_{n,q}$ be the output weight when the input consists of $n$ independent
Bernoulli variables with parameter $q$, and the wrapping convolution has
memory $\sigma$.  For every $\epsilon>0$, there are constants $q_0>0$ and
$K<\infty$ such that
\begin{equation}
 \Pr[H_{n,q}\leq\lfloor\delta n\rfloor]
 \leq K\exp\left(
  -nq\bigl(I^{\mathrm{w}}_{\sigma,\delta}-\epsilon\bigr)
 \right)
 \label{eq:wrapped-ec-uniform-sparse-tail}
\end{equation}
for every $n\geq1$ and $q\in(0,q_0]$.  For every $q_*>0$, there are constants
$K_*,j_*>0$ such that
\begin{equation}
 \Pr[H_{n,q}\leq\lfloor\delta n\rfloor]
 \leq K_*e^{-j_*n}
 \label{eq:wrapped-ec-uniform-intermediate-tail}
\end{equation}
for every $n\geq1$ and $q\in[q_*,1/2]$.
\end{lemma}

\begin{proof}
The probability generating function is
\begin{equation}
 \Ebb[z^{H_{n,q}}]
 =\mathbf e_\sigma^\mathsf T
  \mathbf N_{\sigma,q}(z)^n\mathbf 1.
 \label{eq:wrapped-ec-output-pgf}
\end{equation}
For fixed $q>0$ and $z>0$, the matrix is irreducible.  Perron--Frobenius
theory therefore bounds \eqref{eq:wrapped-ec-output-pgf} by a constant times
$\rho_\sigma(q,z)^n$.  Chernoff extraction then gives the exponent in
\eqref{eq:wrapped-ec-tail-rate}.

We first identify the sparse limit.  At $(q,z)=(0,1)$, the state space has two
closed classes.  The active class consists of states $0,\ldots,\sigma-1$, and
the zero class consists of state $\sigma$.  Define
\begin{equation}
 c_\sigma:=\frac{1}{2^\sigma-1},
 \qquad
 \alpha_\sigma:=\frac{2^{\sigma-1}}{2^\sigma-1}.
 \label{eq:wrapped-ec-class-parameters}
\end{equation}
Within the active class, the stationary probability of state $\ell$ is
$\alpha_\sigma2^{-\ell}$.  Hence $c_\sigma$ is the stationary probability of
state $\sigma-1$, and $\alpha_\sigma$ is the mean output weight per step.

Fix $\theta\geq0$ and set $z=e^{-\theta q}$.  First-order perturbation on the
two-dimensional $1$-eigenspace gives the following rates.  The active class
enters the zero class at rate $c_\sigma$, and the zero class enters the active
class at rate one.  The marker adds the killing rate
$\theta\alpha_\sigma$ within the active class.  The resulting effective matrix
is
\begin{equation}
 G_\sigma(\theta):=
 \begin{pmatrix}
  -c_\sigma-\theta\alpha_\sigma & c_\sigma\\
  1 & -1
 \end{pmatrix}.
 \label{eq:wrapped-ec-effective-generator}
\end{equation}
Indeed, let the right basis consist of the indicators of the two closed
classes, and let the left basis consist of their stationary distributions.
Restricting the first derivative of
$\mathbf N_{\sigma,q}(e^{-\theta q})$ to these dual bases gives
\eqref{eq:wrapped-ec-effective-generator}.
Its rows index the active and zero classes.  The larger eigenvalue is
\begin{equation}
 g_\sigma(\theta)
 =-\frac{1+c_\sigma+\theta\alpha_\sigma}{2}
 +\frac{
  \sqrt{(1-c_\sigma-\theta\alpha_\sigma)^2+4c_\sigma}
 }{2}.
 \label{eq:wrapped-ec-effective-eigenvalue}
\end{equation}
Consequently,
\begin{equation}
 \ln\rho_\sigma(q,e^{-\theta q})
 =qg_\sigma(\theta)+O(q^2).
 \label{eq:wrapped-ec-sparse-spectral-expansion}
\end{equation}
Substitution in \eqref{eq:wrapped-ec-tail-rate} yields
\begin{equation}
 \liminf_{q\downarrow0}
 \frac{J^{\mathrm{w}}_{\sigma,\delta}(q)}{q}
 \geq
 \sup_{\theta\geq0}
 \{-\delta\theta-g_\sigma(\theta)\}.
 \label{eq:wrapped-ec-effective-optimization}
\end{equation}
Set $x:=\delta/\alpha_\sigma$.  Since $\delta<1/2$, the supremum equals
\begin{equation}
 \left(\sqrt{1-x}-\sqrt{c_\sigma x}\right)^2
 =I^{\mathrm{w}}_{\sigma,\delta}.
 \label{eq:wrapped-ec-effective-rate}
\end{equation}
This proves \eqref{eq:wrapped-ec-sparse-exponent-limit}.

The maximizing $\theta$ in \eqref{eq:wrapped-ec-effective-optimization} is
positive.  The Perron vector for
$\mathbf N_{\sigma,q}(e^{-\theta q})$ converges to the positive Perron vector
of \eqref{eq:wrapped-ec-effective-generator}, lifted constantly over each
closed class.  Its largest and smallest coordinates therefore have a bounded
ratio for all sufficiently small $q$.  Applying the componentwise
Perron-vector bound to \eqref{eq:wrapped-ec-output-pgf}, and then using
\eqref{eq:wrapped-ec-sparse-spectral-expansion}, proves
\eqref{eq:wrapped-ec-uniform-sparse-tail}.

It remains to treat parameters bounded away from zero.  At $z=1$,
$\mathbf N_{\sigma,q}(1)$ is stochastic.  Its stationary probabilities obey
\begin{equation}
 \pi_\ell=\pi_0 2^{-\ell}\quad(1\leq\ell<\sigma),
 \qquad
 \pi_\sigma=\pi_{\sigma-1},
 \qquad
 \pi_0=\frac12.
 \label{eq:wrapped-ec-stationary-law}
\end{equation}
The derivative of $\ln\rho_\sigma(q,z)$ with respect to $\ln z$ at $z=1$ is
therefore $1/2$.  Since $\delta<1/2$, a marker below one makes the objective
in \eqref{eq:wrapped-ec-tail-rate} positive.  Continuity follows from
Perron--Frobenius theory and compact restriction of the marker.  A finite
cover of $q\in[q_*,1/2]$ then proves
\eqref{eq:wrapped-ec-uniform-intermediate-tail}.

Finally, every row of $\mathbf N_{\sigma,1/2}(z)$ sums to $(1+z)/2$.
Thus $\rho_\sigma(1/2,z)=(1+z)/2$.  Optimizing over $z$ gives
\eqref{eq:wrapped-ec-dense-exponent}.
\end{proof}

\begin{theorem}[Binary wrapped Expand--Convolute distance]
\label{thm:wrapped-ec-linear-distance}
Fix an integer $\sigma\geq1$.  Let $k=k(n)$ satisfy $k/n\to R$ for a constant
$R\in(0,1)$.  Sample $\B\gets\dist^{\mathrm{Ber}}_{k,n,p_n}$ with
\begin{equation}
 p_n=C\frac{\ln n}{n}
 \label{eq:wrapped-ec-density}
\end{equation}
for a constant $C>0$.  Independently sample a wrapping convolution $\C$ of
memory $\sigma$.  Fix $\delta\in(0,1/2)$.  If
\begin{align}
 C I^{\mathrm{w}}_{\sigma,\delta}&>1,
 \label{eq:wrapped-ec-sparse-condition}\\
 \mathsf{h}(\delta)&<(1-R)\ln2,
 \label{eq:wrapped-ec-gv-condition}
\end{align}
then
\begin{equation}
 \Pr_{\B,\C}\left[
  \min_{\xv\in\Fbb_2^k\setminus\{0\}}
  \hw{\xv\B\C}\leq\lfloor\delta n\rfloor
 \right]=o(1).
 \label{eq:wrapped-ec-linear-distance-conclusion}
\end{equation}
In particular, with probability $1-o(1)$, the sampled generator is injective
and its code has relative minimum distance greater than $\delta$.
\end{theorem}

\begin{proof}
Write $q_r=(1-(1-2p_n)^r)/2$.  We bound the exact first moment in
\eqref{eq:wrapped-expand-convolute-first-moment} in three message-weight
ranges.

First consider sparse messages.  Lemma~\ref{lem:wrapped-ec-endpoint-exponents}
and \eqref{eq:wrapped-ec-sparse-condition} give constants
$\epsilon>0$ and $a_0>0$ such that
\begin{equation}
 \gamma:=
 C\bigl(I^{\mathrm{w}}_{\sigma,\delta}-\epsilon\bigr)(1-a_0)>1.
 \label{eq:wrapped-ec-sparse-margin}
\end{equation}
We choose $a_0$ small enough that
\eqref{eq:wrapped-ec-uniform-sparse-tail} applies whenever $rp_n\leq a_0$.
The bounds
\begin{equation}
 q_r\leq rp_n,
 \qquad
 q_r\geq rp_n(1-a_0)
 \label{eq:wrapped-ec-qr-sparse-bounds}
\end{equation}
then give
\begin{equation}
 \sum_{r\leq a_0/p_n}
 \binom{k}{r}
 \Pr[H_{n,q_r}\leq\lfloor\delta n\rfloor]
 \leq
 K\sum_{r\geq1}k^r n^{-\gamma r}
 =o(1).
 \label{eq:wrapped-ec-sparse-sum}
\end{equation}

Next consider intermediate weights.  If $rp_n>a_0$, then
\begin{equation}
 q_r\geq q_*:=\frac{1-e^{-2a_0}}{2}>0.
 \label{eq:wrapped-ec-intermediate-q}
\end{equation}
Let $j_*>0$ be the uniform exponent from
\eqref{eq:wrapped-ec-uniform-intermediate-tail}.  Choose
$\eta\in(0,R/2)$ such that
\begin{equation}
 R\mathsf{h}(\eta/R)<j_*.
 \label{eq:wrapped-ec-intermediate-eta}
\end{equation}
The Hamming-ball estimate
\begin{equation}
 \sum_{r\leq\eta n}\binom{k}{r}
 \leq\exp\left(nR\mathsf{h}(\eta/R)+o(n)\right)
 \label{eq:wrapped-ec-intermediate-count}
\end{equation}
shows that the contribution of $a_0/p_n<r\leq\eta n$ is exponentially small.

Finally, suppose that $r>\eta n$.  Uniformly over these weights,
\begin{equation}
 0\leq\frac12-q_r
 \leq\frac12e^{-2p_nr}
 \leq\frac12n^{-2C\eta}
 =o(1).
 \label{eq:wrapped-ec-dense-q}
\end{equation}
Every realized wrapping convolution is an invertible linear map.  The
preimage of the output Hamming ball therefore has at most
$\exp(n\mathsf{h}(\delta)+o(n))$ elements.  Since $q_r\leq1/2$, each input word
has probability at most $(1-q_r)^n$.  Thus every fixed message in this range
has failure probability at most
\begin{equation}
 \exp\left(
  n\bigl(\mathsf{h}(\delta)-\ln2+o(1)\bigr)
 \right).
 \label{eq:wrapped-ec-dense-tail}
\end{equation}
There are at most $2^k$ messages.  Condition
\eqref{eq:wrapped-ec-gv-condition} makes their total contribution $o(1)$.
All three ranges vanish, which proves
\eqref{eq:wrapped-ec-linear-distance-conclusion}.
\end{proof}

For $\sigma=1$, \eqref{eq:wrapped-ec-sparse-rate} equals $I_\delta$ from
\eqref{eq:ea-sparse-rate}.  Thus the theorem recovers the accumulator density
condition.  As $\sigma$ tends to infinity,
$I^{\mathrm{w}}_{\sigma,\delta}$ tends to $1-2\delta$.  At
$\delta=0.05$ and $\sigma=21$, condition
\eqref{eq:wrapped-ec-sparse-condition} becomes
\begin{equation}
 C>1.111623.
 \label{eq:wrapped-ec-delta-five-density}
\end{equation}
For $n=5\mathbin{\cdot}2^{20}$, this threshold corresponds to expected
Bernoulli row weight greater than $17.20$.  The dense condition remains the
binary Gilbert--Varshamov condition.

For comparison, the nonwrapping theorem of
\cite[Theorem~3]{EC} requires memory proportional to $\log n$ and imposes
\begin{equation}
 C>\frac{2}{(20/41-\delta)^2},
 \qquad
 R<
 \frac{e-1}{(e+1)\ln2}(20/41-\delta)^2.
 \label{eq:wrapped-ec-prior-asymptotic-conditions}
\end{equation}
At $\delta=0.05$, these conditions require $C>10.4344$ and $R<0.1278$.
The new wrapping theorem instead requires $C>1.111623$ for memory $21$ and
permits every $R<0.7136$.  The latter rate is the binary
Gilbert--Varshamov threshold at this distance.

The fixed-memory wrapping tables in \cite[Appendix~A.1]{EC} use a comparison
with the nonwrapping chain.  The paper states that its small-memory calculation
assumes the conjectural extension of its Lemma~4.  Theorem
\ref{thm:wrapped-ec-linear-distance} analyzes the wrapping transfer matrix
directly.  It therefore gives an unconditional fixed-memory statement.

The same transfer matrix composes with the exact-row expander without an
interleaver.  Define the wrapping slice tail
\begin{equation}
 Q^{\textsf{wConv}}_{n,L}(u)
 :=\frac{1}{\binom{n}{u}}
   \sum_{h=0}^{L}[X^uY^h]\,
   \mathbf e_\sigma^\mathsf{T}
   \mathbf W_\sigma(X,Y)^n\mathbf 1.
 \label{eq:wrapping-slice-tail}
\end{equation}
Lemma~\ref{lem:exchangeable-serial-composition} and the positive shell
recurrence in \eqref{eq:row-shell-recurrence} give
\begin{equation}
 \boxed{
 \mu_{\mathrm{wEC}}^{\mathrm{row}}(L)
 =\sum_{r=1}^{k}\binom{k}{r}
   \sum_{u=0}^{n}\rho_{r,u}
   Q^{\textsf{wConv}}_{n,L}(u).
 }
 \label{eq:row-wrapped-expand-convolute-first-moment}
\end{equation}
This identity is exact and contains only nonnegative terms.  For positive
markers $a$ and $z\leq1$, its slice tail satisfies
\begin{equation}
 Q^{\textsf{wConv}}_{n,L}(u)
 \leq
 \frac{a^{-u}z^{-L}}{\binom{n}{u}}
 \mathbf e_\sigma^\mathsf{T}
 \mathbf W_\sigma(a,z)^n\mathbf 1.
 \label{eq:wrapping-slice-chernoff}
\end{equation}
The input marker extracts the shell, and the output marker extracts the lower
tail.  This two-marker bound is the direct interface between the fixed-row
expander and the wrapped convolution.

The region-stratified left-regular ensemble has a different symmetry: it is
exchangeable within each region, not across all $n$ coordinates.  The
convolution state must therefore be retained between regions.  Define the
matrix-valued slice enumerator
\begin{equation}
 \mathbf S_{\ell,u}(Y)
 :=\frac{1}{\binom{\ell}{u}}[X^u]\,
      \mathbf W_\sigma(X,Y)^\ell
 \label{eq:wrapping-region-slice-matrix}
\end{equation}
and, for a message of weight $r$, define
\begin{equation}
 \mathbf R_{\ell,r}(Y)
 :=\sum_{u=0}^{\ell}\rho^{(\ell)}_{r,u}
      \mathbf S_{\ell,u}(Y).
 \label{eq:wrapping-regular-region-matrix}
\end{equation}

\begin{theorem}[Exact left-regular wrapped-EC first moment]
\label{thm:wrapping-regular-first-moment}
Assume that $n=d\ell$.  Sample
$\B\gets\dist^{\mathrm{reg}}_{k,n,d}$ and an independent wrapping convolution
$\C$ of memory $\sigma$, and set $\G:=\B\C$.  Then the expected number of
nonzero messages whose encoding has weight at most $L$ is
\begin{equation}
 \boxed{
 \mu_{\mathrm{wEC}}^{\mathrm{reg}}(L)
 =\sum_{r=1}^{k}\binom{k}{r}\sum_{h=0}^{L}[Y^h]\,
   \mathbf e_\sigma^\mathsf{T}
   \mathbf R_{\ell,r}(Y)^d\mathbf 1.
 }
 \label{eq:wrapping-regular-first-moment}
\end{equation}
Consequently,
\begin{equation}
 \Pr[\exists\xv\ne0:\ \hw{\xv\G}\leq L]
 \leq\mu_{\mathrm{wEC}}^{\mathrm{reg}}(L).
 \label{eq:wrapping-regular-failure}
\end{equation}
\end{theorem}

\begin{proof}
Fix a message of weight $r$.  Lemma
\ref{lem:region-left-regular-shell-law} makes its input in one region a mixture
of uniform Hamming slices.  Lemma~\ref{lem:wrapping-transfer-enumerator}, with
specified initial and final states, shows that its averaged transfer matrix is
\eqref{eq:wrapping-regular-region-matrix}.  Inputs and convolution taps are
independent across regions.  Matrix multiplication retains the convolution
state at each regional boundary, so the $d$ regions give
$\mathbf R_{\ell,r}(Y)^d$.  Extract the lower output-weight tail, sum over
messages, and apply Markov's inequality.
\end{proof}

The exact matrix coefficient in
\eqref{eq:wrapping-regular-region-matrix} is practical for small message
weights.  Lemma~\ref{lem:regular-poisson-comparison} gives a simpler
all-weight bound.

\begin{lemma}[Poissonized left-regular transfer bound]
\label{lem:wrapping-regular-poisson-bound}
For every $r\geq1$ and $z\in(0,1]$,
\begin{equation}
 \Ebb\!\left[z^{\hw{\xv\B\C}}\right]
 \leq
 \gamma_r^d\,
 \mathbf e_\sigma^\mathsf{T}
 \mathbf N_{\sigma,q^{\mathrm{reg}}_r}(z)^n\mathbf 1
 \label{eq:wrapping-regular-poisson-mgf}
\end{equation}
for each fixed weight-$r$ message $\xv$.  Hence
\begin{equation}
 \Pr[\exists\xv\ne0:\ \hw{\xv\B\C}\leq L]
 \leq
 \sum_{r=1}^{k}\binom{k}{r}\gamma_r^d z^{-L}
 \mathbf e_\sigma^\mathsf{T}
 \mathbf N_{\sigma,q^{\mathrm{reg}}_r}(z)^n\mathbf 1.
 \label{eq:wrapping-regular-poisson-first-moment}
\end{equation}
\end{lemma}

\begin{proof}
Average $z^{\hw{\yv\C}}$ over the independent convolution to obtain a
nonnegative function of $\yv$.  Apply Lemma~\ref{lem:regular-poisson-comparison}
to this function, then use the Bernoulli-input transfer matrix from
\eqref{eq:wrapping-message-active}--\eqref{eq:wrapping-message-zero}.  The
lower-tail moment bound and a union bound over messages give
\eqref{eq:wrapping-regular-poisson-first-moment}.
\end{proof}

Stirling's formula gives $\gamma_r=\Theta(\sqrt r)$.  Thus the comparison
penalty is polynomial in $r$ for fixed left degree $d$.  In contrast, obtaining
an exact-weight row by conditioning every Bernoulli row separately incurs a
constant penalty once for each selected message row.  The regional
construction is therefore especially useful in the intermediate-weight part
of a distance proof.

For dense messages, a direct point-mass bound avoids the Poisson conditioning
penalty.  Let $V_n(L):=\sum_{h=0}^{L}\binom{n}{h}$.

\begin{lemma}[Dense left-regular point-mass bound]
\label{lem:wrapping-regular-dense-bound}
Assume that $n=d\ell$.  For every fixed weight-$r$ message and every invertible
binary linear map $\C$,
\begin{equation}
 \Pr[\hw{\xv\B\C}\leq L]
 \leq
 V_n(L)\,2^{d-n}
 \left(1+e^{-2r/\ell}\right)^n.
 \label{eq:wrapping-regular-dense-tail}
\end{equation}
Consequently, the weight-$r$ first-moment contribution is at most
\begin{equation}
 \binom{k}{r}V_n(L)\,2^{d-n}
 \left(1+e^{-2r/\ell}\right)^n.
 \label{eq:wrapping-regular-dense-first-moment}
\end{equation}
\end{lemma}

\begin{proof}
Lemma~\ref{lem:regular-point-mass} bounds every point of the expander image.
Invertibility makes the preimage under $\C$ of an output Hamming ball have
cardinality $V_n(L)$.  Multiplying the point-mass bound by this cardinality proves
\eqref{eq:wrapping-regular-dense-tail}; multiplying by $\binom{k}{r}$ proves
\eqref{eq:wrapping-regular-dense-first-moment}.
\end{proof}

The factor $2^d$ records the fixed parity of every region.  The remaining
mixing overhead is at most
$n\log_2(1+e^{-2r/\ell})$ bits.  At rate one half and message weight $k/2$,
this overhead becomes
\begin{equation}
 d+n\log_2(1+e^{-d/2}).
 \label{eq:wrapping-regular-rate-half-overhead}
\end{equation}
Thus constant left degree can give a positive relative distance, but a family
that approaches the GV benchmark requires growing degree.  At finite length,
\eqref{eq:wrapping-regular-rate-half-overhead} gives a direct way to select the
degree before running the exact transfer certificate.

The Poissonized bound can cover a contiguous message-weight block with one
transfer calculation.  For $1\leq s\leq t\leq k/2$, the functions
$q^{\mathrm{reg}}_r$ and $\gamma_r$ are nondecreasing.  The latter claim
follows from
\[
 \frac{\gamma_{r+1}}{\gamma_r}
 =e\left(\frac{r}{r+1}\right)^r>1.
\]
Moreover,
$\mathbf N_{\sigma,q}(z)=(1-q)\mathbf W_\sigma(q/(1-q),z)$, and
$\mathbf W_\sigma(a,z)^n$ has nonnegative coefficients in $a$.  Therefore,
for every $r\in[s,t]$ and $z\in(0,1]$, the right-hand side of
\eqref{eq:wrapping-regular-poisson-mgf} is at most
\begin{equation}
 \gamma_t^d(1-q^{\mathrm{reg}}_s)^n
 \mathbf e_\sigma^\mathsf{T}
 \mathbf W_\sigma\!\left(
  \frac{q^{\mathrm{reg}}_t}{1-q^{\mathrm{reg}}_t},z
 \right)^n\mathbf 1.
 \label{eq:wrapping-regular-uniform-block}
\end{equation}
Multiplication by $z^{-L}\sum_{r=s}^t\binom{k}{r}$ gives a valid
first-moment bound for the whole block.

\begin{theorem}[Certified rate-$1/2$ left-regular bound]
\label{thm:wrapped-ec-certified-regular-rate-half}
Set
\[
 k:=1{,}048{,}572,
 \quad d:=28,
 \quad \ell:=74{,}898,
 \quad n:=d\ell=2k,
 \quad L:=230{,}730.
\]
Sample $\B\gets\dist^{\mathrm{reg}}_{k,n,d}$ and an independent zero-state
wrapping convolution $\C$ of memory $9$.  Then
\begin{equation}
 \Pr_{\B,\C}[\exists\xv\ne0:\ \hw{\xv\B\C}\leq L]
 <2.036627\mathbin{\cdot}10^{-7}<2^{-22.2273}.
 \label{eq:wrapped-ec-certified-regular-rate-half}
\end{equation}
The certified relative distance is
\[
 \frac{L}{n}=0.1100210572\ldots
 =0.9999381317\ldots\,\delta_{\mathrm{GV}}(1/2).
\]
\end{theorem}

\begin{proof}
The verifier \path{scripts/regular_ec_certificate.py} evaluates message
weights $1$ through $10$ using the normalized regional slice matrices in
\eqref{eq:wrapping-region-slice-matrix}.  Their total contribution is below
$1.113280\mathbin{\cdot}10^{-35}$.  It covers weights $11$ through
$k/4-1$ with $23$ Poissonized transfer blocks from Lemma
\ref{lem:wrapping-regular-poisson-bound}; their total contribution is below
$3.146936\mathbin{\cdot}10^{-12}$.  Finally, it applies Lemma
\ref{lem:wrapping-regular-dense-bound} in $3072$ positive blocks of width at
most $256$.  Their contribution is below
$2.036595\mathbin{\cdot}10^{-7}$.  The verifier parses every selected marker
as a fixed decimal and evaluates all three regimes with outward-rounded Arb
balls at $192$-bit precision.  Summing them proves
\eqref{eq:wrapped-ec-certified-regular-rate-half}.
\end{proof}

At essentially the same length, the memory-$9$ Bernoulli profile in Theorem
\ref{thm:wrapped-ec-certified-rate-half-gv} has expected row weight $42.2127$
and cutoff $230{,}741$.  The left-regular profile replaces that random row
weight by exactly $28$, a reduction of $33.67\%$, while losing $11$ output
weight units.  Its distance remains within $14.276$ output-weight units of the
binary GV shell.  This comparison measures expander sparsity; an encoder-cost
comparison must also account for the regional sampling and transpose.

The intermediate-weight certificate uses one transfer bound for a contiguous
range of message weights.  Set $p:=d/n$ and let $q_r$ be defined by
\eqref{eq:bernoulli-expander-parity}.  Write
\begin{equation}
 F_\sigma(a,z):=
 \mathbf e_\sigma^\mathsf{T}\mathbf W_\sigma(a,z)^n\mathbf 1.
 \label{eq:wrapping-bivariate-transform}
\end{equation}

\begin{lemma}[Uniform wrapped-EC block bound]
\label{lem:wrapping-uniform-block}
Fix integers $1\leq\ell\leq h\leq k/2$ and $z\in(0,1]$.  For every
$r\in[\ell,h]$, a Bernoulli expander of density $p$ followed by an independent
wrapping convolution satisfies
\begin{equation}
 \Pr[\hw{\xv\B\C}\leq L]
 \leq z^{-L}(1-q_\ell)^n
 F_\sigma\!\left(\frac{q_h}{1-q_h},z\right)
 \label{eq:wrapping-uniform-block-tail}
\end{equation}
for each fixed message $\xv$ of weight $r$.  Consequently, the exact-row
first-moment contribution of all such messages is at most
\begin{equation}
 \vartheta_{n,d}^{-h}
 \left(\sum_{r=\ell}^{h}\binom{k}{r}\right)
 z^{-L}(1-q_\ell)^n
 F_\sigma\!\left(\frac{q_h}{1-q_h},z\right),
 \label{eq:wrapping-uniform-exact-row-block}
\end{equation}
where $\vartheta_{n,d}$ is defined in
\eqref{eq:exact-row-conditioning-probability}.
\end{lemma}

\begin{proof}
Fix $r\in[\ell,h]$ and set $a_r:=q_r/(1-q_r)$.  Averaging the input marker in
Lemma~\ref{lem:wrapping-transfer-enumerator} over independent Bernoulli
coordinates gives
\[
 \Ebb[z^{\hw{\xv\B\C}}]=(1-q_r)^nF_\sigma(a_r,z).
\]
The sequence $q_r$ is nondecreasing.  The polynomial $F_\sigma(a,z)$ has
nonnegative coefficients in $a$.  Hence
\[
 (1-q_r)^nF_\sigma(a_r,z)
 \leq(1-q_\ell)^nF_\sigma(a_h,z).
\]
The lower-tail moment bound proves
\eqref{eq:wrapping-uniform-block-tail}.  Apply
Lemma~\ref{lem:exact-row-conditioning} to each message.  Its conditioning
penalty is at most $\vartheta_{n,d}^{-h}$.  Summing over the block proves
\eqref{eq:wrapping-uniform-exact-row-block}.
\end{proof}

\begin{theorem}[Certified exact-row wrapped-EC bound]
\label{thm:wrapped-ec-certified-exact-row}
Let $(s,d,U,b)$ be any row of the following table, where $b$ is a certified
lower bound on $-\log_2 U$.
\[
\begin{array}{c|c|c|c}
s&d&U&b\\ \hline
20&15&5.160307473\mathbin{\cdot}10^{-7}&20.886039\\
25&17&4.259825504\mathbin{\cdot}10^{-7}&21.162702\\
30&20&4.830790936\mathbin{\cdot}10^{-8}&24.303165
\end{array}
\]
Set $k:=2^s$, $n:=5k$, $L:=n/20$, and $\sigma:=21$.  Independently sample
\[
 \B\gets\dist^{\mathrm{row}}_{k,n,d}
 \quad\text{and a wrapping convolution }\C\text{ of memory }\sigma.
\]
Then
\begin{equation}
 \Pr_{\B,\C}\left[
   \exists\xv\neq0:\ \hw{\xv\B\C}\leq L
 \right]
 \leq U<2^{-b}<2^{-20}.
 \label{eq:wrapped-ec-certified-exact-row-result}
\end{equation}
\end{theorem}

\begin{proof}
For $s=20,25,30$, respectively, the Arb verifier evaluates message weights
through $59$, $240$, and $529$ with the positive shell recurrence and
\eqref{eq:wrapping-slice-chernoff}.  It checks $879$, $4073$, and $5291$
reachable convolution shells.  The resulting small-weight contributions are
at most
\[
 4.315211437\mathbin{\cdot}10^{-7},\qquad
 1.671977463\mathbin{\cdot}10^{-7},\qquad
 7.040348863\mathbin{\cdot}10^{-9}.
\]

The verifier partitions the remaining weights below $k/4$ into $51$, $86$,
and $119$ contiguous blocks.  Each block uses a fixed decimal marker in
\eqref{eq:wrapping-uniform-exact-row-block}.  The three intermediate-weight
contributions are at most
\[
 8.450960361\mathbin{\cdot}10^{-8},\qquad
 2.587848041\mathbin{\cdot}10^{-7},\qquad
 4.126756049\mathbin{\cdot}10^{-8}.
\]
The dense L2 bound from Lemma~\ref{lem:exact-row-dense} applies from $k/4$
onward.  Every realized wrapping convolution is invertible, so the preimage of
the output Hamming ball has the same cardinality.  The three dense
contributions are below $2^{-60{,}355}$, $2^{-6{,}001{,}120}$, and
$2^{-350{,}267{,}009}$.

The verifier parses all markers as fixed decimals.  It evaluates matrix
powers, shell probabilities, conditioning probabilities, and the final sum
with outward-rounded Arb balls at $192$-bit precision.  Summing the three
regimes gives the displayed values.  Equation
\eqref{eq:row-wrapped-expand-convolute-first-moment} and Markov's inequality
imply \eqref{eq:wrapped-ec-certified-exact-row-result}.
\end{proof}

\subsection{Finite comparison with the prior EC analysis}
\label{sec:wrapped-ec-finite-comparison}

We first hold the code ensemble fixed and change only the analysis.  Set
$k:=2^{20}$, $n:=5k$, $L:=n/20$, and $\sigma:=5$.  For
$C_0\in\{3,2.5,2.3\}$, set
\begin{equation}
 p:=C_0\frac{\ln n}{n}
 \label{eq:wrapped-ec-comparison-density}
\end{equation}
and sample $\B\gets\dist^{\mathrm{Ber}}_{k,n,p}$.  Independently sample a
wrapping convolution $\C$ of memory $\sigma$, initialized in the zero state.
These parameters match the nonsystematic, rate-$1/5$, zero-state rows of
\cite[Table~4]{EC}.

\begin{theorem}[Certified Bernoulli wrapped-EC comparison]
\label{thm:wrapped-ec-certified-bernoulli-comparison}
For each row below, the exact-transfer column upper-bounds
\[
 \Pr_{\B,\C}\left[
  \exists\xv\neq0:\ \hw{\xv\B\C}\leq L
 \right].
\]
The final column is a certified lower bound on the corresponding number of
security bits.
\[
\begin{array}{c|c|c|c|c}
C_0&\Ebb[\text{row weight}]
 &\text{prior bound}&\text{exact transfer}&\text{bits}\\ \hline
3.0&46.4171&2.26\mathbin{\cdot}10^{-7}
 &2.187816572\mathbin{\cdot}10^{-10}&32.089789\\
2.5&38.6810&1.72\mathbin{\cdot}10^{-5}
 &9.485867880\mathbin{\cdot}10^{-8}&23.329645\\
2.3&35.5865&1.73\mathbin{\cdot}10^{-4}
 &1.124903039\mathbin{\cdot}10^{-6}&19.761768
\end{array}
\]
\end{theorem}

\begin{proof}
The Arb verifier evaluates message weights $1$ through $32$ individually with
\eqref{eq:wrapped-expand-convolute-chernoff}.  For
$C_0=3,2.5,2.3$, the resulting contributions are at most
\[
 1.954065032\mathbin{\cdot}10^{-10},\qquad
 9.482273889\mathbin{\cdot}10^{-8},\qquad
 1.124788945\mathbin{\cdot}10^{-6}.
\]
The verifier covers the remaining weights below $k/4$ with $9$, $10$, and
$11$ uniform transfer blocks.  Lemma~\ref{lem:wrapping-uniform-block} gives
intermediate contributions at most
\[
 2.337515400\mathbin{\cdot}10^{-11},\qquad
 3.593990580\mathbin{\cdot}10^{-11},\qquad
 1.140938107\mathbin{\cdot}10^{-10}.
\]
From $k/4$ onward, invertibility bounds each preimage by the output Hamming
ball.  The dense contribution is below $10^{-740{,}000}$ in every row.

The verifier parses every marker and each value of $C_0$ as a fixed decimal.
It evaluates all matrix powers, activation probabilities, and final sums with
outward-rounded Arb balls at $192$-bit precision.  The exact first moment in
\eqref{eq:wrapped-expand-convolute-first-moment} and Markov's inequality give
the displayed bounds.
\end{proof}

The exact transfer improves the three Table~4 bounds by factors of at least
$1032$, $181$, and $153$.  Equivalently, it gains $10.01$, $7.50$, and $7.26$
certified bits.  Table~4 explicitly assumes the conjectural extension of
Lemma~4 in \cite{EC}; the new calculation does not use that comparison lemma.

Tables~5 and~6 of \cite{EC} initialize the chain with its stationary
distribution.  That boundary condition changes the sampled construction, so
those rows are not included in the same-ensemble comparison.  The exact-row
certificate in Theorem~\ref{thm:wrapped-ec-certified-exact-row} is also a
different design point.  At $k=2^{20}$, it replaces expected Bernoulli row
weight $46.4171$ and memory $5$ by exact row weight $15$ and memory $21$.
Its failure bound remains below $2^{-20}$.  This comparison moves work from
the expander to the convolution; it is not by itself an encoding-time
comparison.

\subsection{Finite approach to the Gilbert--Varshamov distance}
\label{sec:wrapped-ec-finite-gv}

The asymptotic dense-weight condition in
Theorem~\ref{thm:wrapped-ec-linear-distance} is the binary
Gilbert--Varshamov condition.  A finite certificate must also retain the
subexponential factors and the residual bias of the expander output.  The
following bound keeps both effects without evaluating additional transfer
matrices.

For integers $0\leq t<N/2$, define the binary Hamming-ball volume
\[
 V_N(t):=\sum_{j=0}^{t}\binom Nj.
\]

\begin{lemma}[Dense blocks for wrapped EC]
\label{lem:wrapped-ec-dense-block}
Fix $1\leq\ell\leq h\leq k$ and let
\[
 \mathcal X_{\ell,h}
 :=\{\xv\in\{0,1\}^k:\ell\leq\hw{\xv}\leq h\}.
\]
For the Bernoulli wrapped-EC ensemble of
Theorem~\ref{thm:wrapped-expand-convolute-first-moment},
\begin{equation}
 \sum_{\xv\in\mathcal X_{\ell,h}}
 \Pr_{\B,\C}[\hw{\xv\B\C}\leq L]
 \leq
 |\mathcal X_{\ell,h}|V_n(L)(1-q_\ell)^n.
 \label{eq:wrapped-ec-dense-block}
\end{equation}
Moreover,
\begin{equation}
 V_N(t)
 \leq
 \binom Nt\frac{N-t+1}{N-2t+1}.
 \label{eq:hamming-ball-geometric-bound}
\end{equation}
\end{lemma}

\begin{proof}
Fix $\xv\in\mathcal X_{\ell,h}$ with weight $r$.  The coordinates of
$\xv\B$ are independent Bernoulli variables with parameter $q_r$.
Since $q_r\leq1/2$, every vector has probability at most $(1-q_r)^n$.
Every realized wrapping convolution is invertible.  Conditional on $\C$, the
preimage of the radius-$L$ output ball therefore has $V_n(L)$ elements.
The sequence $q_r$ is nondecreasing, so the probability for $\xv$ is at most
$V_n(L)(1-q_\ell)^n$.  Summing over $\mathcal X_{\ell,h}$ proves
\eqref{eq:wrapped-ec-dense-block}.

For $j<t$, consecutive binomial coefficients satisfy
\[
 \frac{\binom Nj}{\binom N{j+1}}
 =\frac{j+1}{N-j}
 \leq\frac{t}{N-t+1}.
\]
The geometric-series bound gives
\eqref{eq:hamming-ball-geometric-bound}.
\end{proof}

Let $\delta_{\mathrm{GV}}(R)$ be the solution in $(0,1/2)$ to
\begin{equation}
 h(\delta_{\mathrm{GV}}(R))=(1-R)\ln2.
 \label{eq:binary-gv-distance}
\end{equation}
At rate $R=1/5$, this value is
$\delta_{\mathrm{GV}}(1/5)=0.2430038538\ldots$.

\begin{theorem}[Certified finite GV comparison]
\label{thm:wrapped-ec-certified-finite-gv}
Set $k:=2^{20}$ and $n:=5k$.  For each row below, set
$p:=C_0\ln(n)/n$ and use a zero-state wrapping convolution of memory
$\sigma$.  Independently sample
\[
 \B\gets\dist^{\mathrm{Ber}}_{k,n,p}
 \quad\text{and}\quad \C.
\]
The last column gives a certified value $b$ such that
\begin{equation}
 \Pr_{\B,\C}\left[
  \exists\xv\neq0:\ \hw{\xv\B\C}\leq L
 \right]<2^{-b}.
 \label{eq:wrapped-ec-finite-gv-result}
\end{equation}
\[
\begin{array}{c|c|c|c|c|c}
C_0&\sigma&C_0\ln n&L/n
 &(L/n)/\delta_{\mathrm{GV}}(1/5)&b\\ \hline
3&5 &46.4171&0.119999886&49.3819\%&20.104148\\
3&13&46.4171&0.159999847&65.8425\%&20.161525\\
4&5 &61.8895&0.189999962&78.1880\%&20.280537\\
4&9 &61.8895&0.233999825&96.2947\%&20.166272\\
4&13&61.8895&0.241799927&99.5046\%&20.030269\\
5&9 &77.3619&0.242617607&99.8411\%&29.140327\\
5&9 &77.3619&0.243002319&99.9994\%&21.320038
\end{array}
\]
\end{theorem}

\begin{proof}
The six integer cutoffs are
\[
 629145,\quad 838860,\quad 996147,\quad
 1226833,\quad 1267728,\quad 1272015,\quad 1274032.
\]
The verifier evaluates message weights $1$ through $32$ with the exact
transfer matrix.  It covers weights $33$ through $k/4-1$ with, respectively,
$15$, $14$, $21$, $22$, $21$, $21$, and $21$ uniform transfer blocks.

For the dense range, the verifier applies
Lemma~\ref{lem:wrapped-ec-dense-block} to ten blocks.  The block endpoints are
\[
 \frac{k}{4},\frac{3k}{10},\frac{7k}{20},\frac{2k}{5},
 \frac{9k}{20},\frac{47k}{100},\frac{12k}{25},
 \frac{49k}{100},\frac{99k}{200},\frac{k}{2},k+1,
\]
with integer floors where necessary.  Equation
\eqref{eq:hamming-ball-geometric-bound} bounds both the message counts and the
output ball.  The first six dense contributions are below $10^{-1002}$.
The final dense contribution is at most
$3.800902566\mathbin{\cdot}10^{-7}$.

The resulting all-weight bounds, in table order, are
\[
\begin{split}
 &8.872546060\mathbin{\cdot}10^{-7},\quad
  8.526605175\mathbin{\cdot}10^{-7},\quad
  7.851452089\mathbin{\cdot}10^{-7},\\
 &8.498595923\mathbin{\cdot}10^{-7},\quad
  9.338732260\mathbin{\cdot}10^{-7},\\
 &1.690003011\mathbin{\cdot}10^{-9},\quad
  3.819697244\mathbin{\cdot}10^{-7}.
\end{split}
\]
The verifier parses each marker as a fixed decimal and uses outward-rounded
Arb balls at $192$-bit precision.  The first-moment bound in
Theorem~\ref{thm:wrapped-expand-convolute-first-moment} proves
\eqref{eq:wrapped-ec-finite-gv-result}.
\end{proof}

The last row is within $0.000001535$ of the binary GV distance.  Equivalently,
its integer cutoff is $8.0451$ below $n\delta_{\mathrm{GV}}(1/5)$.  It reaches
$99.99937\%$ of the asymptotic benchmark with expected expander row weight
$77.3619$ and convolutional memory $9$.  The preceding row uses the same code
ensemble and retains $29.140327$ bits at a slightly smaller distance.  The GV
value is not a finite-length upper bound; the comparison measures the
remaining gap to the standard asymptotic random-linear-code benchmark.

\subsubsection{Rate one half.}

At rate $1/2$, the binary GV distance is
\[
 \delta_{\mathrm{GV}}(1/2)=0.1100278644\ldots.
\]
The next theorem gives four cost profiles at this rate.  The table orders the
profiles by memory; the expected expander row weight falls as memory grows.

\begin{theorem}[Certified rate-$1/2$ GV comparison]
\label{thm:wrapped-ec-certified-rate-half-gv}
Set $k:=2^{20}$, $n:=2k$, and $L:=230741$.  For each row below, set
$p:=C_0\ln(n)/n$ and use a zero-state wrapping convolution of memory
$\sigma$.  Independently sample
\[
 \B\gets\dist^{\mathrm{Ber}}_{k,n,p}
 \quad\text{and}\quad \C.
\]
Then
\[
 \Pr_{\B,\C}\left[
  \exists\xv\neq0:\ \hw{\xv\B\C}\leq L
 \right]\leq U<2^{-b}.
\]
\[
\begin{array}{c|c|c|c|c}
C_0&\sigma&C_0\ln n&U&b\\ \hline
3.15&5&45.8517&7.637727109\mathbin{\cdot}10^{-7}&20.320353\\
3&7&43.6683&5.278295360\mathbin{\cdot}10^{-7}&20.853424\\
2.9&9&42.2127&5.560776634\mathbin{\cdot}10^{-7}&20.778210\\
2.85&13&41.4849&5.971032143\mathbin{\cdot}10^{-7}&20.675516
\end{array}
\]
Both rows certify relative distance
\[
 \frac{L}{n}=0.1100258827\ldots
 =0.9999819889\ldots\,\delta_{\mathrm{GV}}(1/2).
\]
\end{theorem}

\begin{proof}
The verifier evaluates message weights $1$ through $32$ with the exact
transfer matrix.  In table order, these contributions are at most
\[
\begin{split}
 &4.248414142\mathbin{\cdot}10^{-7},\quad
  1.884923212\mathbin{\cdot}10^{-7},\\
 &2.164751311\mathbin{\cdot}10^{-7},\quad
  2.574039056\mathbin{\cdot}10^{-7}.
\end{split}
\]
The four transfer regimes end at weights $26362$, $26362$, $27461$, and
$28045$.  They use $11$, $10$, $9$, and $9$ uniform blocks.  Their
contributions are at most
\[
\begin{split}
 &8.192194653\mathbin{\cdot}10^{-11},\quad
  3.128007837\mathbin{\cdot}10^{-10},\\
 &2.989483906\mathbin{\cdot}10^{-10},\quad
  1.589442947\mathbin{\cdot}10^{-10}.
\end{split}
\]

After each transfer regime, the verifier applies
Lemma~\ref{lem:wrapped-ec-dense-block}.  Its positive point-mass partitions
contain $12$, $17$, $17$, and $17$ blocks.  Equation
\eqref{eq:hamming-ball-geometric-bound} gives dense contributions at most
\[
\begin{split}
 &3.388493748\mathbin{\cdot}10^{-7},\quad
  3.390244141\mathbin{\cdot}10^{-7},\\
 &3.393035839\mathbin{\cdot}10^{-7},\quad
  3.395403644\mathbin{\cdot}10^{-7}.
\end{split}
\]
The verifier checks all transfer markers, block endpoints, Hamming-ball
bounds, and final sums with outward-rounded Arb balls at $192$-bit precision.
Summing the three regimes proves the stated bounds.
\end{proof}

The GV shell lies at output weight $230745.156\ldots$.  Thus the certified
integer cutoff is $4.156$ output-weight units below the asymptotic benchmark.
At the next integer cutoff, the current dense bound already exceeds the
$2^{-20}$ target.  Further progress therefore requires a sharper dense
partition or a different finite-length argument.

\subsection{Legacy combinatorial derivation (draft)}

The remainder of this section records the earlier zone-based derivation for
comparison.  Lemma~\ref{lem:wrapping-transfer-enumerator} supersedes it as the
proof-facing enumerator.

Consider the following example with $\sigma = 4$:
\begin{align*}
  \xv=&\texttt{ 0\textcolor{blue}{1}\red{0101000}\green{0}\red{101}\textcolor{blue}{1}00\textcolor{blue}{1}\red{00100100101}}\\
  \yv=&\texttt{ 0\olive{1001}\green{1000}\textcolor{blue}{10000}00\olive{11001110}\textcolor{violet}{1000}}\\
  \text{ on?}=&\texttt{ 0111111111111000111111111111}%\\
  %v=&\texttt{ 000000001000000000000000}\\
  %t_0=&\texttt{ 001100000000000000110001}
\end{align*}

Compared to the prior section, we define the zones of the pattern differently:
\begin{itemize}

  \item \textbf{Deadzones:} As before, deadzones are runs of zeros and occur at the very beginning or after terminations. Deadzones are aligned in $\xv$ and $\yv$ and indicate when the convolution is off. Deadzones may be empty.

  \item \textbf{Terminations:} This is the sequence $10^{\sigma}$ in $\yv$ and must have a \blue{1} in $\xv$ that is aligned with the last \blue{0} in $\yv$. For example, $\sigma=4$ we have
    \begin{align*}
      \xv=& \texttt{?????\blue{1}?}\\
      \yv=&\texttt{?\blue{10000}?}
    \end{align*}
    After this, the convolution transitions into a deadzone consisting of $\texttt{0}^*$.

  \item \textbf{Near-terminations:} This is the sequence of $\green{10}^{\sigma-1}$ in $\yv$ with the constraint that it must be followed by a $\texttt{1}$ and this training \texttt{1} in $\yv$ must be aligned with a $\green{0}$ in $\xv$. For example, if $\sigma=4$ we must have
    \begin{align*}
      \xv=& \texttt{?????\green{0}?}\\
      \yv=&\texttt{?\green{1000}1?}
    \end{align*}
    We do not count the trailing \texttt{1} in $\yv$ as part of the termination pattern but a constraint on what follows.

  \item \textbf{Input Freezones:} These are sequences of \texttt{\red{0}}s and \texttt{\red{1}}s in $\xv$ when the convolution is on. In $x$, freezones start after the first  \texttt{\textcolor{blue}{1}} and end in \texttt{\textcolor{blue}{1}}.

  \item \textbf{Output Freezones:} In $y$, freezones can contain up to $\sigma-2$ consecutive \texttt{\olive{0}}s as they would otherwise turn into either a termination or a near-termination.

  \item \textbf{Trailzone:} This is the sequence $\violet{10}^{\sigma-1}$ at the very end of $\yv$. It is similar to a near termination but since it is at the end, it does not result in the corresponding requirement of having a \green{0} in $\xv$. As a result, we must make it a special case.
\end{itemize}

\subsection{Peter's Wrapping}

\paragraph{Additonal Constraints.}
We will add the following constraints to the pattern:
\begin{itemize}
  \item There are exactly $r$ run transitions, where the convolution turns on or off. I.e., $r:=\left|\{i\in [n] \mid \qv_{i}\neq \qv_{i-1} \}\right|$ and $\qv_0:=1,\qv_i:=1$ if $\yv_{i-[\sigma]}\neq0^\sigma$ and 0 otherwise.
  \item There are exactly $t_0$ output freezone zeros. I.e. $t_0:=\left|\{i \in [n] \mid \yv_i=0,\qv_{i+[\sigma]}=1^\sigma,\}\right|$
  \item There are exactly $v$ near-termination zones. I.e. $v:=|\{ i \mid \yv_{i+[\sigma+1] = \texttt{\green{10}}^{\sigma}\texttt{\green{1}}}\}|$
  \item There are exactly $d\in\{0,1\}$ trailzones. Note $r$ being odd implies $d=0$.
\end{itemize}

Given the values of the constraints, we will then enumerate all $(\xv,\yv,\C)$ that are compatible.

\paragraph{Useful quantities.}
Let us first define useful quantities that are implied by our constraints.
\begin{itemize}
  \item $r_0:=\lfloor r/2\rfloor$ terminations $\blue{10}^\sigma$ in $\yv$. Equivalently, there are $r_0$ runs of deadzones not counting the first one.
  \item $r_1:=\lceil r/2\rceil$ non-zero runs.
  \item $z:=n-h-(v+d)(\sigma -1)-r_0\sigma -t_0$  deadzones \texttt{0}s. There are $n-h$ zeros in general, $(v+d)(\sigma -1)$ of them are used for near termiations and trailing zones, $r_0\sigma$ used for terminations, and $t_0$ are used in the output freezone.
  \item $f_0:=n-w-z-v$ input freezone \red{0}s . There are $n-w$ in general, $z$ used in deadzones and $v$ used in near terminations.
  \item $f_1:=w-r$ input freezone \red{1}s. There are $w$ ones in general and each of the $r$ run transitions must occure at a \blue{1} in $\xv$.
  \item $t_1:=h-v-r_0-d$ output freezone $\olive{1}$s. $h$ in general, $v$ are used by near terminations, $r_0$ are used for terminations, and $d$ are used for the trailing zone.
\end{itemize}

\paragraph{Baseline.}

Let us first count place the $t_1$  output freezone $\olive{1}$s  and the $d$ trailzones $\violet{10}^{\sigma-1}$. There is exactly one way to do this. Our example above with $t_1=7,d=1$ would look like
\begin{align*}
  x&=\texttt{?\blue{1}???????????????}\\
  y&=\texttt{?\olive{1}?\olive{1}?\olive{1}?\olive{1}?\olive{1}?\olive{1}?\olive{1}?\textcolor{violet}{1000}}
\end{align*}

\paragraph{Near terminations.} Near-terminations can occur before any of the $t_1$ output freezone \texttt{\olive{1}}s or one of the two special cases at the end: (1) before the $d \in \{0,1\}$ trailzones when we finish in an on state, or (2) before the last termination when we finish in a deadzone, i.e. if $1-r_1+r_0=1$. For 2, we only count the position before the last termination as it is not equivalent to being before an output freezone \olive{1}, unlike the other terminations. Each of these configurations results in a different $y$. There are
$$
E_1=\mathcal{B}(v,t_1+d+(1-r_1+r_0))
$$
different ways of ordering the \olive{1}s, \textcolor{violet}{1000}s and near-terminations. In our example we have
\begin{align*}
  x&=\texttt{?\blue{1}???????\green{0}???????????????}\\
  y&=\texttt{?\olive{1}?\olive{1}?\green{1000}??\olive{1}?\olive{1}?\olive{1}?\olive{1}?\olive{1}?\textcolor{violet}{1000}}
\end{align*}
Note that a zero cannot follow a near termination in $\yv$.

\paragraph{Terminations.} There are exactly $r_0$ terminations in general. However, if the convolution ends in the off state, then clearly the last termination must come right before the last deadzone. Therefore, there are $r_1-1$ terminations that are free to be placed, one for each of the $r_1$ times the convolution is on minus the last (because the last run either does not terminate or the last termination is fixed).
A termination may come before each of the $t_1$ output freezone \olive{1}s or $v$ near terminations $\green{10}^{\sigma-1}$, or at the very end if (1) there is a trailzone, i.e. $d=1$, or (2) the convolution is in the off state at the end, i.e. $1-r_1+r_0=1$. Therefore, there are
$$
E_2=\ballsBins(r_1-1, t_1+v+d+(1-r_1+r_0))
$$
such different ways of ordering the terminations for each ordering of the \texttt{\olive{1}}s and near-terminations. In our example, we have
\begin{align*}
  x&=\texttt{?\blue{1}???????\green{0}???\blue{1}?\blue{1}?????????????}\\
  y&=\texttt{?\olive{1}?\olive{1}?\green{1000}\blue{10000}?\olive{1}?\olive{1}?\olive{1}?\olive{1}?\olive{1}?\textcolor{violet}{1000}}
\end{align*}
All \texttt{1}s in the output are now placed and therefore in our example we can eliminate the \texttt{?} after the near termination, as we know a zero can not follow it. Additionally, the start of each non-zero run in $\xv$ is now known and is depicted above with a \blue{1}.

\paragraph{Deadzones.}
The $r_0+1$ deadzones can only occur at the beginning or after terminations and each way results in a different $y$. There are $z$ zeros in the deadzones. Therefore, there are
$$
E_3=\mathcal{B}(z,r_0+1)
$$
such different ways of ordering the \texttt{0}s in the deadzones for each ordering of the \texttt{\olive{1}}s, near-terminations, and terminations. In our example we have
\begin{align*}
  x&=\texttt{0\blue{1}???????\green{0}???\blue{1}00\blue{1}?????????????}\\
  y&=\texttt{0\olive{1}?\olive{1}?\green{1000}\blue{10000}00\olive{1}?\olive{1}?\olive{1}?\olive{1}?\olive{1}?\textcolor{violet}{1000}}
\end{align*}

\paragraph{Freezone output zeros.}
The $t_0$ output freezone {\olive{0}}s can only come after the $t_1$ output freezone {\olive{1}}s and there can be up to $\sigma-2$ consecutive \olive{0}s at a time. Therefore, there are
$$
E_4=N(t_0,t_1,\sigma-2)
$$ such different ways of ordering the \texttt{\olive{0}}s in the freezones for each ordering of the \texttt{\olive{1}}s, near-terminations, terminations, and \texttt{0}s. In our example we have
\begin{align*}
  x&=\texttt{0\blue{1}???????\green{0}???\blue{1}00\blue{1}???????????}\\
  y&=\texttt{0\olive{1001}\green{1000}\blue{10000}00\olive{11001110}\textcolor{violet}{1000}}
\end{align*}

\paragraph{Freezone inputs.}
What remains in the ordering of the $f_0+f_1$ input freezone \texttt{\red{0}}s and \texttt{\red{1}}s. Each ordering results in a different $x$. Therefore, there are
$$
E_5=\binom{f_0+f_1}{f_0}
$$
such different ways of ordering the \texttt{\red{0}}s of $x$ for each. At this stage, we have also completely determined $x$. In our example we have our final input output string
\begin{align*}
  x&=\texttt{0\textcolor{blue}{1}\red{010100}\green{0}\red{0101}\textcolor{blue}{1}00\textcolor{blue}{1}\red{00100100101}}\\
  y&=\texttt{0\olive{1001}\green{1000}\blue{10000}00\olive{11001110}\textcolor{violet}{1000}}\\
  G&=\texttt{0011111101111000011111111111}
\end{align*}

\paragraph{Putting it together.}
Thus, for fixed values of $n,w,h,r,d,v,t_0$, the number of possible consistent $x$ and $y$ is
\begin{align*}
  \enum_{w,h,r,d,v,t_0} &= 2^{-g}E_1E_2E_3E_4E_5 P_{w,h,r,d, v, t_0}
  %&=2^{-g}\binom{v}{t_1+1} \binom{r'-b}{h-v-r'+b} \binom{z}{r'+1} \binom{f_0+f_1}{f_0} \\
  %&\quad \left(\sum_{i=0}^{t_1+v}(-1)^{i}\binom{t_1+v}{i}\binom{t_0+t_1+v-i(\sigma-1)-1}{t_1+v-1}\right)
\end{align*}
where $g:=f_1+f_0$ is the number of times we assumed the result of the convolution times the non-zero state is equal to some value, i.e. the positions where $G_i=1$ in the example above. $P$ is the bounds check predicate which is defined as
$$
P_{w,h,r,d, v, t_0} :=
\begin{cases}
  0 & \text{if } z<0 \vee f_0 <0 \vee f_1 <0 \vee t1 < 0 \vee h-r_0-d < 0 \vee\\
  &\quad  k - r - v - f_0 - f_1 < 0 \vee r_1-r_0 -d < 0\\
  1 & \text{otherwise }
\end{cases}
$$
Unlike the other enumerator, the bounds check is infact required due to prevent intermidiate counts from going negative. A reparameterization may be possible to avoid this issue.

To obtain the enumerator for fixed values of $n,w,h$, we thus need to sum over all possible values of $r,d,v,t_0$.
Thus, for fixed values of $n,w,h$, the number of possible consistent $x$ and $y$ is
$$
\enum_{w,h} = \sum_{r\geq 0}\sum_{d\in \{0, 1\}}\sum_{v\geq 0}\sum_{t_0\geq 0}\mathbb{E}_{n,w,h,r,d,v,t_0}
$$

\subsection{Srini's Wrapping}

\noindent Let us first describe what $y$ looks like in general:
\begin{itemize}
  \item $y$ may begin with a deadzone.
  \item $y$ may then contain a mixture of freezones, near-terminations, terminations, and deadzones, subject to the following restrictions:
    \begin{itemize}
      \item Deadzones only follow terminations.
      \item Freezones cannot follow freezones. Also, freezones following a termination or a deadzone must begin with a \texttt{\red{1}}.
    \end{itemize}
  \item $y$ may end with a trailzone.
\end{itemize}
Based on $y$, $x$ then contains:
\begin{itemize}
  \item deadzones corresponding to deadzones in $y$
  \item freezones corresponding to terminations, near-terminations, freezones and a trailzone in $y$, subject to the following restrictions:
    \begin{itemize}
      \item freezones corresponding to terminations in $y$ end in \texttt{\textcolor{blue}{1}}
      \item freezones corresponding to near-terminations in $y$ end in \texttt{\green{0}}
      \item freezones corresponding to any sequence following a termination or a deadzone in $y$ begin with \texttt{\textcolor{blue}{1}}
    \end{itemize}
\end{itemize}

\noindent We will have the independent variables $n,w,h,r,b,d,v$ which represent:
\begin{itemize}
  \item the lengths of $x$ and $y$, $n=28$
  \item the weight of the input $x$, $w=11$
  \item the weight of the output $y$, $h=10$
  \item the number of times the convolution turns on, $r=2$
  \item the indicator of whether the convolution turned off at the end, $b=0$
  \item the indicator of whether the convolution ends in a trailzone, $d=1$
  \item the number of near-terminations, $v=1$
  \item the number of \texttt{\red{0}}s in freezones in $y$, $t_0=5$
\end{itemize}
Using these independent parameters, we can define several others as below:
\begin{itemize}
  \item the number of terminations, $r'=r+b-1=1$
  \item the number of \texttt{0}s in deadzones, $z=n-h-(v+d)(\sigma-1)-r'\sigma-t_0=3$
  \item the number of \texttt{\red{0}}s in freezones in $x$, $f_0=n-w-z-v=13$
  \item the number of \texttt{\red{1}}s in freezones in $x$, $f_1=w-r-r'=8$
  \item the number of \texttt{\red{1}}s in freezones in $y$, $t_1=h-2v-r'-d=6$
\end{itemize}
All of the above can be argued and summarized using the table below:
\begin{center}
  \begin{tabular}{|c|c|c|}\hline
    \textbf{Property} & \textbf{Input $x$} & \textbf{Output $y$}  \\\hline
    \texttt{\red{0}}s & $f_0$ & $t_0$ \\\hline
    \texttt{\textcolor{blue}{0}}s & $-$ & $r'\sigma$ \\\hline
    \texttt{\green{0}}s & $v$ & $v(\sigma-1)$ \\\hline
    \texttt{\textcolor{violet}{0}}s & $-$ & $d(\sigma-1)$ \\\hline
    \texttt{0}s & $z$ & $z$ \\\hline\hline
    \texttt{\red{1}}s & $f_1$ & $t_1$ \\\hline
    \texttt{\textcolor{blue}{1}}s & $r+r'$ & $r'$ \\\hline
    \texttt{\green{1}}s & $-$ & $2v$ \\\hline
    \texttt{\textcolor{violet}{1}}s & $-$ & $d$ \\\hline
    \texttt{1}s & $-$ & $-$ \\\hline
  \end{tabular}
\end{center}

\noindent Suppose we fix the values of $n,w,h,r,b,d,v,t_0$. We will now describe how to enumerate over all possible $y$ and $x$ consistent with these values. Firstly, if $d=1$, we place a trailzone at the end of $y$. Conceptually, we will build up the pattern and count the number of options. At this stage, we have placed all \texttt{\textcolor{violet}{0}}s and \texttt{\textcolor{violet}{1}}s of the tailzone in $y$ for which there is exactly one way to do it. With respect to our exmaple, at this point we have
\begin{align*}
  x&=\texttt{?????}\\
  y&=\texttt{?\textcolor{violet}{1000}}
\end{align*}
where \texttt{?} denodes an undertermined string.

\paragraph{Freezone output ones.}
Next, we consider the freezone \texttt{\olive{1}}s in $y$. Given the current containts, these will all appear before the tailzone and there is one way to place them, e.g.
\begin{align*}
  x&=\texttt{?\blue{1}???????????????}\\
  y&=\texttt{?\olive{1}?\olive{1}?\olive{1}?\olive{1}?\olive{1}?\olive{1}?\textcolor{violet}{1000}}
\end{align*}

\paragraph{Near terminations.} Near-terminations can occur arbitrarily with respect to the \texttt{\olive{1}}s and each way results in a different $y$. There are
$$
E_1=\mathcal{B}(v,t_1+1)
$$
such different ways of ordering the \texttt{\olive{1}}s and near-terminations. At this stage, we have also placed all output freezone \texttt{\olive{1}}s, near termination \texttt{\green{0}}s and \texttt{\green{1}}s and trailzone \texttt{\textcolor{violet}{0}}s and \texttt{\textcolor{violet}{1}}s in $y$. In our example we have
\begin{align*}
  x&=\texttt{?\blue{1}?????\green{0}???????????????}\\
  y&=\texttt{?\olive{1}?\green{10001}?\olive{1}?\olive{1}?\olive{1}?\olive{1}?\olive{1}?\textcolor{violet}{1000}}
\end{align*}

\paragraph{Output terminations.}
Next, consider any fixed ordering of the \texttt{\olive{1}}s and near-terminations in $y$. Terminations can occur \emph{almost} arbitrarily with respect to the \texttt{\olive{1}}s and near-terminations and each way results in a different $y$. The only restriction has to do with $b$. By definition, $b=1$ if and only if there is a termination after the last \texttt{\olive{1}} or near-termination and $d = 0$. Therefore, there are
$$
E_2=\mathcal{B}(r'-b,t_1+v+b+d)=\mathcal{B}(r'-b,h-v-r'+b)
$$
such different ways of ordering the terminations for each ordering of the \texttt{\olive{1}}s and near-terminations.
At this stage, we have also placed all \texttt{\textcolor{blue}{0}}s and \texttt{\textcolor{blue}{1}}s in $y$. In out example we have
\begin{align*}
  x&=\texttt{?\blue{1}?????\green{0}?????\blue{1}?\blue{1}?????????????}\\
  y&=\texttt{?\red{1}?\green{10001}?\blue{10000}?\red{1}?\red{1}?\red{1}?\red{1}?\red{1}?\textcolor{violet}{1000}}
\end{align*}

\paragraph{Deadzones.}
Next, consider any fixed ordering of the \texttt{\red{1}}s, near-terminations, and terminations in $y$. Deadzones can only occur at the beginning or after terminations and each way results in a different $y$. Therefore, there are
$$
E_3=\mathcal{B}(z,r'+1)
$$
such different ways of ordering the \texttt{0}s in the deadzones for each ordering of the \texttt{\red{1}}s, near-terminations, and terminations. At this stage, we have also placed all \texttt{0}s in $y$. In our example we have
\begin{align*}
  x&=\texttt{0\blue{1}?????\green{0}?????\blue{1}00??????????????}\\
  y&=\texttt{0\red{1}?\green{10001}?\blue{10000}00\red{1}?\red{1}?\red{1}?\red{1}?\red{1}?\textcolor{violet}{1000}}
\end{align*}

\paragraph{Freezone output zeros.}
Next, consider any fixed ordering of the \texttt{\red{1}}s, near-terminations, terminations, and \texttt{0}s in $y$. The \texttt{\red{0}}s can only come after \texttt{\red{1}}s or near-terminations and there can be up to $\sigma-2$ of them consecutively and each way results in a different $y$. Therefore, there are$$E_4=N(t_0,t_1+v,\sigma-2)$$such different ways of ordering the \texttt{\red{0}}s in the freezones for each ordering of the \texttt{\red{1}}s, near-terminations, terminations, and \texttt{0}s. At this stage, we have also placed all \texttt{\red{0}}s in $y$ and hence completely determined it. In our example we have
\begin{align*}
  x&=\texttt{0\blue{1}??????\green{0}????\blue{1}00????????????}\\
  y&=\texttt{0\red{100}\green{10001}\blue{10000}00\red{11001110}\textcolor{violet}{1000}}
\end{align*}

\paragraph{Freezone inputs.}
Next, consider any fixed $y$. Corresponding to $y$, we can place all the \texttt{\green{0}}s, \texttt{0}s, and \texttt{\textcolor{blue}{1}}s in $x$. The remaining spots can be filled in arbitrarily with \texttt{\red{0}}s and \texttt{\red{1}}s and each way results in a different $x$. Therefore, there are
$$
E_5=\binom{f_0+f_1}{f_0}
$$
such different ways of ordering the \texttt{\red{0}}s of $x$ for each $y$. At this stage, we have also completely determined $x$. In our example we have our final input output string
\begin{align*}
  x&=\texttt{0\textcolor{blue}{1}\red{010100}\green{0}\red{0101}\textcolor{blue}{1}00\textcolor{blue}{1}\red{00100100101}}\\
  y&=\texttt{0\red{100}\green{10001}\blue{10000}00\red{11001110}\textcolor{violet}{1000}}\\
  G&=\texttt{0011111101111000011111111111}
\end{align*}

\paragraph{Putting it together.}
Thus, for fixed values of $n,w,h,r,b,d,v,t_0$, the number of possible consistent $x$ and $y$ is
\begin{align*}
  \mathbb{E}_{w,h,r,b,d,v,t_0} &= 2^{-g}E_1E_2E_3E_4E_5\\
  &=2^{-g}\binom{v}{t_1+1} \binom{r'-b}{h-v-r'+b} \binom{z}{r'+1} \binom{f_0+f_1}{f_0} \\
  &\quad \left(\sum_{i=0}^{t_1+v}(-1)^{i}\binom{t_1+v}{i}\binom{t_0+t_1+v-i(\sigma-1)-1}{t_1+v-1}\right)
\end{align*}
where $g:=n-z-v-2r-1+b$ is the number of times we assumed the result of the convolution times the non-zero state is equal to some value, i.e. the positions where $G_i=1$ in the example above.
To obtain the enumerator for fixed values of $n,w,h$, we thus need to sum over all possible values of $r,b,d,v,t_0$. Note that:
\begin{itemize}
  \item $r$ can take values from $0$ through $\min\{w,h,\frac{n-h}{\sigma}+1\}$.
  \item By definition, $b=1\implies d=0$, therefore, $(b, d) \in \{(0,0), (0,1), (1,0)\}$.
  \item $v$ can take values from $0$ through $\min\{n-w,\frac{h}{2},\frac{n-h}{\sigma-1}\}$.
  \item $t_0$ can take values from $0$ through $n-h$.
\end{itemize}
Thus, for fixed values of $n,w,h$, the number of possible consistent $x$ and $y$ is$$\mathbb{E}_{w,h} = \sum_{r=0}^{\min\{w,h,\frac{n-h}{\sigma}+1\}}\sum_{b, d\in \{0, 1\}:bd\ne 1}\sum_{v=0}^{\min\{n-w,\frac{h}{2},\frac{n-h}{\sigma-1}\}}\sum_{t_0=0}^{n-h}\mathbb{E}_{n,w,h,r,b,d,v,t_0}$$

\subsection{Extended Convolution}\label{sec:extendedconv}

Let $\mathbf{x} \in \{0,1\}^k$ be an input vector, and let $\mathbf{C}$ be a fixed convolution matrix, inducing the linear map $\mathbf{y} = \mathbf{x}\mathbf{C}$.
We observe a structural issue when $\mathbf{x}$ is low-weight and only non-zero near its tail.

Consider the extreme example $\mathbf{x} = 000\ldots0001$.
Because the convolution operator receives only a single nonzero input located at the final coordinate, the output $\mathbf{y} = \mathbf{x}\mathbf{C}$ has nonzero support only near the end of the vector.  A subsequent permutation of $\mathbf{y}$ may, with non-negligible probability, place this nonzero entry again among the final coordinates.  In such a case, the composition of convolution and permutation fails to achieve meaningful diffusion, leaving us effectively where we started.

More generally, the same issue arises when only the last $O(1)$ positions of $\mathbf{x}$ are nonzero, e.g., $\mathbf{x} = 000\ldots01111$.
Then $\mathbf{y} = \mathbf{x}\mathbf{C}$ still exhibits support concentrated near the final coordinates $\mathbf{y} = 000\ldots01???$, and after permutation the support may again lie in the last few positions.
This degeneracy does not occur when the nonzero entries of $\mathbf{x}$ reside in the middle or beginning of the vector, in which case the convolution naturally spreads support over a larger portion of the output.

\medskip
\noindent\textbf{Mitigating the issue.}
We remedy this vulnerability by introducing a so-called \emph{extended convolution}.
In extended convolution, we increase the convolution horizon beyond the nominal length $k$ by producing $\kappa$ additional output positions.
This forces the nonzero tail of $\mathbf{x}$ to mix into a region subsequently permuted into substantially different coordinates.


Consider the input that has trailing zeros (consider the extension part of the input x)

\paragraph{Baseline.}

Let us first count place the $t_1$  output freezone $\olive{1}$s  and the $d$ trailzones $\violet{10}^{\sigma-1}$. There is exactly one way to do this. Our example above with $t_1=7,d=1$ would look like
\begin{align*}
  x&=\texttt{?\blue{1}???????????????000}\\
  y&=\texttt{?\olive{1}?\olive{1}?\olive{1}?\olive{1}?\olive{1}?\olive{1}?\olive{1}?\textcolor{violet}{1000}}
\end{align*}
